\documentclass[11pt,letterpaper]{article}
\usepackage[margin=.68in,top=.58in,bottom=.62in]{geometry}
% Shared style for Saange Tamang's academic and outreach CVs.
% TeX Gyre Pagella ships with standard TeX distributions; no proprietary fonts.
\usepackage[T1]{fontenc}
\usepackage{tgpagella}
\usepackage{xcolor,enumitem,tabularx,microtype,fancyhdr}
\definecolor{ink}{HTML}{202124}
\definecolor{accent}{HTML}{4F685A}
\definecolor{muted}{HTML}{56545A}
\definecolor{rulegray}{HTML}{C9C5C6}
\usepackage[colorlinks=true,urlcolor=accent,linkcolor=accent,
  pdfauthor={Saange Tamang},pdfsubject={Wireless Communications and Machine Learning},
  pdfkeywords={wireless communications, machine learning, reinforcement learning, graph learning}]{hyperref}
\color{ink}
\setlength{\parindent}{0pt}
\setlength{\parskip}{2pt}
\setlength{\emergencystretch}{2em}
\setlist[itemize]{leftmargin=13pt,label=\textbullet,itemsep=3pt,topsep=3pt,parsep=0pt,partopsep=0pt}
\newcommand{\sect}[1]{\par\vspace{8pt}{\color{accent}\fontsize{11}{13}\selectfont\bfseries #1}\par\vspace{2pt}{\color{rulegray}\hrule height .35pt}\vspace{4pt}}
\newcommand{\entry}[2]{\textbf{#1}\hfill{\color{muted}\fontsize{10}{12}\selectfont #2}\par}
\newcommand{\meta}[1]{{\color{muted}\fontsize{10}{12}\selectfont #1}\par}
\newcommand{\paperlink}{\href{https://github.com/Sange-creator/graph-attention-double-dqn-cellular-handover/tree/main/paper}{Manuscript}}
\newcommand{\codelink}{\href{https://github.com/Sange-creator/graph-attention-double-dqn-cellular-handover}{Research code}}
\newcommand{\nameheader}{%
  {\fontsize{28}{31}\selectfont\bfseries Saange Tamang}\par
  \vspace{3pt}
  {\fontsize{10}{12}\selectfont Electronics, Communication \& Information Engineer\enspace |\enspace Licensed Engineer, Nepal Engineering Council}\par
  \vspace{5pt}
  {\fontsize{10}{12}\selectfont Nepal\enspace |\enspace\href{mailto:connectwithsaange@gmail.com}{connectwithsaange@gmail.com}\enspace |\enspace
  \href{https://github.com/Sange-creator}{GitHub}\enspace |\enspace\href{https://www.linkedin.com/in/saange-tamang/}{LinkedIn}}\par
}

\hypersetup{pdftitle={Saange Tamang - Academic CV}}
\pagestyle{fancy}
\fancyhf{}
\renewcommand{\headrulewidth}{0pt}
\fancyfoot[L]{\color{muted}\fontsize{10}{12}\selectfont Saange Tamang\enspace |\enspace Academic CV}
\fancyfoot[R]{\color{muted}\fontsize{10}{12}\selectfont\thepage\ /\ 2}
\begin{document}
\fontsize{10.6}{13.3}\selectfont
\nameheader

\sect{Research interests}
\textbf{Wireless communications:} RAN, 5G/6G, cellular mobility, and network optimization.\par
\textbf{AI/ML:} deep learning, neural networks, graph neural networks, and reinforcement learning.

\sect{Education}
\entry{Paschimanchal Campus, Institute of Engineering}{2022--2026}
B.E. in Electronics, Communication and Information Engineering\par
\meta{Tribhuvan University, Pokhara, Nepal}
\textbf{79.02\% aggregate; First Division.}\par
Selected results: Discrete Structure 95/100; Computer Network 86/100; Data Mining 81/100; Project II 97/100.\par
\vspace{3pt}
\textbf{Higher secondary (+2):} Kathmandu Model College\hfill{\color{muted}\fontsize{10}{12}\selectfont 2019--2021}\par
\textbf{Secondary Education Examination (SEE):} GPA 3.9\hfill{\color{muted}\fontsize{10}{12}\selectfont 2018}

\sect{Conference paper}
\textbf{Saange Tamang}, Sulav Kandel. ``Graph-Attention Double-DQN for Cellular Handover: A Paired Simulation Evaluation.''\par
\meta{19th IOE Graduate Conference, 2026. Accepted for presentation. Proceedings citation pending.}
\paperlink\enspace |\enspace\codelink

\sect{Research experience}
\entry{Graph learning and reinforcement learning for cellular handover}{2026}
\meta{First author: Saange Tamang. Co-author: Sulav Kandel.}
\begin{itemize}
\item Led technical implementation and training. Built the Python simulator and PyTorch controller using a 46-cell graph, 15 features per cell, and action masking.
\item Designed and ran matched comparisons of three frozen policies with a configured LTE A3-style baseline across 639 reserved Kathmandu-area cases and 540 Pokhara transfer groups.
\item Used paired cluster-bootstrap intervals to assess the tradeoff. Kathmandu-area simulations showed 8.86--9.99\% higher aggregate throughput and 36.1--53.0\% fewer handovers per hour, with a higher share of short completed stays.
\end{itemize}
\meta{Sulav Kandel gathered observations that motivated the study and checked site coordinates against OpenStreetMap. We shared manuscript preparation equally.}
\meta{Both evaluations used real road geometry with simulated radio and traffic. Pokhara used synthetic cell layouts. Operator drive-test logs were unavailable for research use.}

\sect{Telecommunications experience}
\entry{Nepal Telecom (NTC), Pokhara}{12 Jan--27 Mar 2026}
\meta{Engineering Intern}
\begin{itemize}
\item Joined a vehicle-based cellular drive test and observed route-based measurement collection. Accompanied a tower maintenance visit and examined installed routers and server-room infrastructure.
\item Explored hybrid mesh topology in Cisco Packet Tracer. Attended training led by senior engineers on cell-site connectivity, DNS, NAT, VPNs, and international upstream links.
\end{itemize}

\newpage
{\fontsize{15}{18}\selectfont\bfseries Saange Tamang}\hfill{\color{muted}\small Academic CV, continued}\par

\sect{Selected engineering projects}
\entry{GNN-DQN-Based Handover Optimization and Load Balancing}{2025--2026}
\textbf{in LTE Networks}\par
\meta{Major project; one year across semesters 7--8; four-person team}
\meta{Supervisor: Asst. Prof. Hom Nath Tiwari}
Collaborated on graph-based RL for cellular handover and load balancing, exploring learned decisions alongside a self-organising-network control layer. Contributed to implementation, evaluation, validation, and analysis; the later conference study narrowed the question to handover optimization.\enspace
\href{https://github.com/Sange-creator/GNN-DQN-Handover-LoadBalancing}{Code/report}

\vspace{7pt}
\entry{SolVote: Electronic Voting and Token Distribution Prototype}{Feb--Aug 2025}
\meta{Undergraduate minor project; semester 6; collaborative engineering project}
Contributed to an academic prototype for registration, token issuance, and voting using React/TypeScript, Node.js/Express, MongoDB, and Solana Web3.js. RFID/fingerprint authentication remained planned integration work.\enspace
\href{https://github.com/Sange-creator/Project-Solvote}{Code}

\vspace{7pt}
\entry{JAMJAMNEPAL: Tourism Application}{Ongoing side project}
Developing a Go/Fiber and MongoDB application for domestic and international tourists visiting Nepal; initial server-side templates and a handler test are in the public repository. AI integration is planned.\enspace
\href{https://github.com/Sange-creator/GO_Tourism_Project}{Code}

\vspace{7pt}
\entry{Computer Network Simulations and Packet Analysis}{Coursework}
Cisco Packet Tracer exercises in IP address configuration, cross-network ping/messages, DHCP, and DNS; Wireshark packet capture/filtering and inspection of DHCP exchanges and DNS queries/responses.\enspace
\href{https://github.com/Sange-creator/Computer-Network}{Lab artifacts}

\sect{Technical preparation}
\begin{tabularx}{\linewidth}{@{}p{1.15in}@{\hspace{9pt}}>{\raggedright\arraybackslash}X@{}}
\textbf{Wireless systems} & LTE Event A3-style handover; RSRP/SINR; mobility and simulated cell-load models; paired service/switching evaluation.\\[5pt]
\textbf{AI/ML methods} & Graph attention; dueling Double DQN; replay buffers; action masking; cluster-bootstrap intervals; transfer-scenario analysis.\\[5pt]
\textbf{Networking labs} & Cisco Packet Tracer, Wireshark, IP address configuration, ICMP/ping, DHCP/DNS simulation, packet capture and filtering.\\[5pt]
\textbf{Research tools} & Python, PyTorch, NumPy, Matplotlib, Git/GitHub, LaTeX; C/C++ (coursework).
\end{tabularx}

\sect{Selected course completions}
\entry{\href{https://www.coursera.org/account/accomplishments/verify/3WMVQTJX2BQC}{Advanced Learning Algorithms}}{12 Sep 2026}
\meta{DeepLearning.AI and Stanford Online, via Coursera}
\vspace{4pt}
\entry{\href{https://www.coursera.org/account/accomplishments/verify/H85L1YRD0MWM}{Supervised Machine Learning: Regression and Classification}}{31 Aug 2026}
\meta{DeepLearning.AI and Stanford Online, via Coursera}
\vspace{4pt}
\entry{\href{https://www.coursera.org/account/accomplishments/verify/P4AO5YNOHQLD}{Claude Code: Software Engineering with Generative AI Agents}}{13 Aug 2026}
\meta{Vanderbilt University, via Coursera}
\vspace{4pt}
\entry{\href{https://verify.skilljar.com/c/3isdp98xhia2}{Introduction to subagents} --- Anthropic Education}{14 Aug 2026}

\sect{Professional qualification}
Licensed Engineer, Nepal Engineering Council.
\end{document}
